\section{Worked problems}
\subsection*{1. The cache boundary}
\textbf{Problem.} A prompt has three tokens. The model has selected its
first continuation but has not evaluated that selection. How many positions
are cached?
\textbf{Solution.} Three, at positions zero through two. Evaluating the
selected token adds position three and advances the count to four.
\subsection*{2. Payload versus allocation}
\textbf{Problem.} Count the default teaching model's KV payload for five
evaluated tokens if stored as packed F32.
\textbf{Solution.} $2\cdot2\cdot1\cdot2\cdot4=32$ bytes per token,
so five tokens require 160 payload bytes. Python object overhead,
allocator metadata and any unused capacity are excluded.
\subsection*{3. A stronger oracle}
\textbf{Problem.} Why is agreement on the greedy token weaker than
agreement on every logit?
\textbf{Solution.} Many different score vectors have the same largest
entry. A faulty cache can change nonwinning scores while preserving that
winner. Compare complete rows at every prefix and include cases with
different token histories and head sharing.
\Needspace{8\baselineskip}
\subsection*{4. Count the actual model}
\textbf{Problem.} Count matrix parameters and hypothetical packed F32
payload for the default decoder, including the untied vocabulary output.
\textbf{Solution.} Each layer has $16+8+8+16+24+24+24=120$ entries.
Embedding and output each add twenty. Total is $20+240+20=280$,
or 1,120 payload bytes at four bytes per entry. Unit normalization gains
and absent biases add no learned parameters.
\Needspace{8\baselineskip}
\subsection*{5. Three predictions, two feedback evaluations}
\textbf{Problem.} Compare projected rows and causal query--key pairs when
predicting three IDs from a three-token prompt, across both layers and
both query heads of the default model.
\textbf{Solution.} Full-prefix lengths are three, four and five: projected
rows total $2(3+4+5)=24$, and pairs total $4(6+10+15)=124$.
Cached evaluation visits five tokens once: ten projected rows and
$4(1+2+3+4+5)=60$ pairs. These are work counts, not timings.
\Needspace{8\baselineskip}
\subsection*{6. Stop without evaluating the end ID}
\textbf{Problem.} Suppose the first selected ID is configured as the end
ID in the trace's policy. How many IDs are selected and how many positions
are cached?
\textbf{Solution.} One ID is selected and retained in the trace. The
three prompt positions remain cached; the selected end ID is not fed back.
Stopping and evaluation are different events.
